\section{Proofs for the bulk}\label{app:bulk}

We prove Theorems~\ref{thm:cp} and~\ref{thm:landau},
Lemma~\ref{lem:landau-mode} and Corollary~\ref{cor:mode-location}. We keep
the notation of Section~\ref{sec:shape} and of Appendix~\ref{app:shape}. So
$\mathcal T$ is the tree of Lemma~\ref{lem:tree}, $\mathcal T_v$ is the family
of $v$, and $\mu_v$ is the mark of $v$. By Lemma~\ref{lem:urn} and the
discussion after Lemma~\ref{lem:tree}, $\lvert\mathcal T_v\rvert$ has the law
$\beta_v$, with mean $N/v$, and $\sum_{v=2}^N\beta_v(d)=N/(d(d+1))$ for
$1\le d\le N-1$. Let
\[
 \Psi_N(\theta)=\sum_{d=1}^{N-1}\frac{e^{i\theta d}-1}{d(d+1)},\qquad
 \Psi_\infty(\theta)=\sum_{d\ge1}\frac{e^{i\theta d}-1}{d(d+1)} .
\]
The law $\CP_N$ has the characteristic function
$\widehat{\CP}_N(\theta)=e^{x\Psi_N(\theta)}$.

\subsection{Proof of Theorem~\ref{thm:cp}}\label{app:bulk-cp}

Here we assume only $N\ge2$ and $0<\varepsilon\le\frac14$. Let
$Z_+=\sum_{v=2}^N\mu_v\lvert\mathcal T_v\rvert$ be the total size of the
flipped families.

\emph{Step 1: $M_N=Z_+$ with high probability.} Two families are nested or
disjoint. Suppose that no marked vertex has a marked strict descendant. Then
the marked families are disjoint, each of their vertices has exactly one
mark on its path to the root, and the other vertices have none. So
$M_N=Z_+$ by Lemma~\ref{lem:tree}. The number of pairs of a marked vertex and
a marked strict descendant has mean
$\varepsilon^2\,\E\sum_v(\lvert\mathcal T_v\rvert-1)$, since marks and tree are
independent. Hence
\[
 \Prob(M_N\ne Z_+)\le\varepsilon^2\sum_{v=2}^N\Bigl(\frac Nv-1\Bigr)
 \le\varepsilon^2N(H_N-1)\le\frac{x^2\log N}{N}.
\]

\emph{Step 2: given the tree, a product.} Fix $\theta\in[-\pi,\pi]$. Given
$\mathcal T$ the marks are independent, so
$\E[e^{i\theta Z_+}\mid\mathcal T]=\prod_{v=2}^N(1+a_v)$ with
$a_v=\varepsilon(e^{i\theta\lvert\mathcal T_v\rvert}-1)$. Here
$\lvert a_v\rvert\le2\varepsilon\le\frac12$ and $\operatorname{Re}a_v\le0$.
Put $W=\sum_va_v$, so that $\operatorname{Re}W\le0$. For complex $a$ with
$\lvert a\rvert\le\frac12$ the series of the logarithm gives
$\lvert\log(1+a)-a\rvert\le\lvert a\rvert^2$. So $\prod_v(1+a_v)=e^{W+\varrho}$
with $\lvert\varrho\rvert\le\sum_v\lvert a_v\rvert^2\le4\varepsilon^2(N-1)$.
Since $\lvert e^w-1\rvert\le\lvert w\rvert e^{\lvert w\rvert}$,
\[
 \Bigl\lvert\prod_v(1+a_v)-e^W\Bigr\rvert\le e^{\operatorname{Re}W}\,
 4\varepsilon^2(N-1)\,e^{4\varepsilon^2(N-1)}\le\frac{4x^2}Ne^{4x^2/N}.
\]

\emph{Step 3: replace $W$ by its mean.} Write $W=\varepsilon F$ with
$F=\sum_{v=2}^N(e^{i\theta\lvert\mathcal T_v\rvert}-1)$. Then
$\E W=\varepsilon\sum_d(e^{i\theta d}-1)\sum_v\beta_v(d)=x\Psi_N(\theta)$,
which is $\log\widehat{\CP}_N(\theta)$. Put $\delta=W-\E W$. Taylor's formula
for $s\mapsto e^{\E W+s\delta}$ on $[0,1]$ gives
\[
 e^W-e^{\E W}-e^{\E W}\delta=\delta^2\int_0^1(1-s)\,e^{\E W+s\delta}\,ds .
\]
The exponent $\E W+s\delta=(1-s)\E W+sW$ has real part at most $0$, so the
right side has modulus at most $\frac12\lvert\delta\rvert^2$. As $\E\delta=0$,
we get $\lvert\E e^W-e^{\E W}\rvert\le\frac12\varepsilon^2\Var F$, where
$\Var F=\E\lvert F-\E F\rvert^2$.

\emph{Step 4: the variance of $F$.} The variable $F$ is a function of the
independent parents $U_1,\dots,U_{N-1}$. Let $F^{(k)}$ be $F$ with $U_k$
replaced by an independent copy $U_k'$. The Efron--Stein
inequality~\cite{efronstein1981} says that
$\Var F\le\frac12\sum_k\E\lvert F-F^{(k)}\rvert^2$. We use it for a complex
function of independent variables that are not identically distributed, so we
recall the proof. Let $\E_k$ be the average over $U_k$ alone, and let
$\zeta_k=\E[F-\E_kF\mid U_1,\dots,U_k]$. Then
$\zeta_k=\E[F\mid U_1,\dots,U_k]-\E[F\mid U_1,\dots,U_{k-1}]$, so
$F-\E F=\sum_k\zeta_k$ and the $\zeta_k$ are orthogonal. Hence
$\Var F=\sum_k\E\lvert\zeta_k\rvert^2\le\sum_k\E\lvert F-\E_kF\rvert^2$, and
$\E\lvert F-\E_kF\rvert^2=\frac12\E\lvert F-F^{(k)}\rvert^2$, because given
the other parents $F$ and $F^{(k)}$ are independent with the same law.

Replacing $U_k$ by $U_k'$ moves the family of $k+1$ from $U_k$ to $U_k'$. So
$\lvert\mathcal T_v\rvert$ changes only for vertices $v\ge2$ on the path from
$U_k$ to the root or on the path from $U_k'$ to the root. These paths do not
depend on the choice at step $k$, since $U_k,U_k'\le k$. Let $\mathrm{dp}(u)$
be the depth of $u$, namely the number of edges between $u$ and the root.
There are at most $\mathrm{dp}(U_k)+\mathrm{dp}(U_k')$ such vertices, and each
term of $F$ moves by at most $2$. Hence
$\E\lvert F-F^{(k)}\rvert^2\le8\,\E[\mathrm{dp}(U_k)^2+\mathrm{dp}(U_k')^2]
=16\,\E\,\mathrm{dp}(U_k)^2$.

Since $\mathrm{dp}(u+1)=1+\mathrm{dp}(U_u)$, the generating function
$g_u(s)=\E s^{\mathrm{dp}(u)}$ satisfies $g_1=1$ and
$u\,g_{u+1}(s)=s\sum_{j\le u}g_j(s)$. Subtracting the same identity for $u-1$
gives $u\,g_{u+1}(s)=(u-1+s)\,g_u(s)$ for $u\ge2$, and $g_2(s)=s$. So
$g_u(s)=\prod_{i=1}^{u-1}\frac{i-1+s}{i}$. Thus $\mathrm{dp}(u)$ is a sum of
independent Bernoulli$(1/i)$ variables with $1\le i<u$, and
$\E\,\mathrm{dp}(u)^2\le H_{u-1}^2+H_{u-1}$. The parent $U_k$ is independent
of the depths of $1,\dots,k$, so $\E\,\mathrm{dp}(U_k)^2\le H_N^2+H_N$. We
get $\Var F\le\frac12(N-1)\cdot16(H_N^2+H_N)\le8N(H_N^2+H_N)$, and the bound
of Step 3 is at most $4x^2(H_N^2+H_N)/N$.

\emph{Step 5: Fourier inversion.} For every integer $m$,
\[
 \Prob(Z_+=m)-\CP_N(m)=\frac1{2\pi}\int_{-\pi}^{\pi}e^{-im\theta}
 \bigl(\E e^{i\theta Z_+}-\widehat{\CP}_N(\theta)\bigr)\,d\theta,
\]
so the left side is at most
$\sup_\theta\lvert\E e^{i\theta Z_+}-\widehat{\CP}_N(\theta)\rvert$ in absolute
value. By Steps 2 to 4 this is at most
$\frac{4x^2}Ne^{4x^2/N}+\frac{4x^2(H_N^2+H_N)}N$. Finally
$\lvert f_N(m)-\Prob(Z_+=m)\rvert\le\Prob(M_N\ne Z_+)$, and Step 1 bounds it.
This proves Theorem~\ref{thm:cp}. \qed

\subsection{Proof of Theorem~\ref{thm:landau}}\label{app:bulk-landau}

We write $B(t)=-it\log\lvert t\rvert-\frac\pi2\lvert t\rvert$ for real
$t\ne0$, and $B(0)=0$. For $t>0$ this is the value of $u\log u$ at $u=-it$,
with the principal logarithm, and $B(-t)=\overline{B(t)}$. We need three
facts about $\Psi_N$ and $\Psi_\infty$.

\emph{Fact 1.} For real $\theta$,
$\lvert\Psi_N(\theta)-\Psi_\infty(\theta)\rvert\le\sum_{d\ge N}\frac2{d(d+1)}=\frac2N$.

\emph{Fact 2.} For $0<\lvert\theta\rvert\le\pi$ we have
$\Psi_\infty(\theta)=-i\theta\log(-i\theta)+R(\theta)$ with
$\lvert R(\theta)\rvert\le\theta^2\bigl(0.631\,\lvert\log\lvert\theta\rvert\rvert+1.53\bigr)$.
The reason is that $\Psi_\infty$ has a closed form in which $\theta$
enters only through $1-e^{i\theta}$, and $1-e^{i\theta}$ is close to
$-i\theta$ for small $\theta$. The term $R$ is what this replacement
changes. We prove the fact in four steps.

We start with the closed form. For $0<\lvert s\rvert\le1$, $s\ne1$,
\[
 \sum_{d\ge1}\frac{s^d-1}{d(d+1)}=\chi(s):=\frac{(1-s)\log(1-s)}s,
\]
because $\frac1{d(d+1)}=\frac1d-\frac1{d+1}$, $\sum_ds^d/d=-\log(1-s)$,
$\sum_ds^d/(d+1)=(-\log(1-s)-s)/s$ and $\sum_d\frac1{d(d+1)}=1$. So
$\Psi_\infty(\theta)=\chi(e^{i\theta})$.

Next we compare $1-e^{i\theta}$ with $-i\theta$. Write
$1-e^{i\theta}=(-i\theta)\,\mathsf u$ with
$\mathsf u=\operatorname{sinc}(\theta/2)\,e^{i\theta/2}$, where
$\operatorname{sinc}y=\sin y/y$ is positive for $\lvert y\rvert\le\frac\pi2$.
So $\mathsf u$ is close to $1$ when $\theta$ is small. The argument of
$-i\theta$ is $\mp\frac\pi2$ for $\pm\theta>0$ and that of $\mathsf u$ is
$\theta/2$. They add up to a number in $(-\frac\pi2,\frac\pi2)$, so
$\log(1-e^{i\theta})=\log(-i\theta)+\log\mathsf u$ for principal logarithms.

Now we get a formula for $R$. By the two steps above,
$\chi(e^{i\theta})=(-i\theta)\,\mathsf ue^{-i\theta}\bigl(\log(-i\theta)+\log\mathsf u\bigr)$.
Subtracting $-i\theta\log(-i\theta)$ gives
\[
 R(\theta)=-i\theta\log(-i\theta)\bigl(\mathsf ue^{-i\theta}-1\bigr)
 -i\theta\,\mathsf ue^{-i\theta}\log\mathsf u .
\]

It remains to bound the three quantities in this formula. We show that
\[
 \lvert\mathsf ue^{-i\theta}-1\rvert\le0.631\,\lvert\theta\rvert,\qquad
 \lvert\log\mathsf u\rvert\le0.53\,\lvert\theta\rvert,\qquad
 \lvert\log(-i\theta)\rvert\le\lvert\log\lvert\theta\rvert\rvert+\frac\pi2 .
\]
For the first, put $y=\theta/2$. Then
$\lvert\mathsf ue^{-i\theta}-1\rvert=\lvert\operatorname{sinc}y-e^{iy}\rvert
\le(1-\operatorname{sinc}y)+\lvert1-e^{iy}\rvert\le\frac{y^2}6+\lvert y\rvert
\le(\frac12+\frac\pi{24})\lvert\theta\rvert$, where the last step uses
$\lvert\theta\rvert\le\pi$, and $\frac12+\frac\pi{24}\le0.631$. For the
second,
$\lvert\log\mathsf u\rvert^2=\log^2\operatorname{sinc}(\theta/2)+\theta^2/4$,
and $-\log\operatorname{sinc}y=\sum_{k\ge1}\frac{\zeta(2k)}{k\pi^{2k}}y^{2k}$
has positive coefficients. So
$\lvert\log\operatorname{sinc}(\theta/2)\rvert/\lvert\theta\rvert$ increases in
$\lvert\theta\rvert$, and its value at $\lvert\theta\rvert=\pi$ gives
$\lvert\log\mathsf u\rvert\le\bigl(\frac14+\frac{\log^2(\pi/2)}{\pi^2}\bigr)^{1/2}\lvert\theta\rvert
\le0.53\,\lvert\theta\rvert$. The third holds because $-i\theta$ has
modulus $\lvert\theta\rvert$ and argument $\pm\frac\pi2$. Since
$\lvert\mathsf ue^{-i\theta}\rvert\le1$, the formula for $R$ now gives
\[
 \lvert R(\theta)\rvert\le\lvert\theta\rvert\Bigl(\lvert\log\lvert\theta\rvert\rvert+\frac\pi2\Bigr)\cdot0.631\,\lvert\theta\rvert
 +\lvert\theta\rvert\cdot0.53\,\lvert\theta\rvert
 \le\theta^2\bigl(0.631\,\lvert\log\lvert\theta\rvert\rvert+1.53\bigr),
\]
as $0.631\cdot\frac\pi2+0.53\le1.53$.

\emph{Fact 3.} For $\lvert\theta\rvert\le\pi$,
$-\operatorname{Re}\Psi_\infty(\theta)=\sum_d\frac{1-\cos\theta d}{d(d+1)}
\ge\frac{\lvert\theta\rvert}{2\pi}$. We may take $\theta\ne0$. Keep only the
terms with $d\le D:=\lfloor\pi/\lvert\theta\rvert\rfloor$ and use
$1-\cos v\ge2v^2/\pi^2$ for $\lvert v\rvert\le\pi$. This gives at least
$\frac{2\theta^2}{\pi^2}\sum_{d\le D}\frac d{d+1}\ge\frac{\theta^2D}{\pi^2}
\ge\frac{\lvert\theta\rvert}{2\pi}$, since $D\ge\frac\pi{2\lvert\theta\rvert}$.

\emph{The function $e^B$ is a characteristic function of class $L$.} Let
$S_x$ have the law $\CP_\infty$ with parameter $x$, namely the law of
$\sum_{d\ge1}d\,\mathcal N_d$. Its characteristic function is
$e^{x\Psi_\infty(\theta)}$. By Fact 2, for $t\ne0$ and
$x\ge\max(1,\lvert t\rvert/\pi)$,
\begin{equation}\label{eq:bulk-AB}
 x\Psi_\infty(t/x)-it\log x=B(t)+xR(t/x),\qquad
 \lvert xR(t/x)\rvert\le\frac{t^2}x\bigl(0.631(\lvert\log\lvert t\rvert\rvert+\log x)+1.53\bigr).
\end{equation}
So the characteristic function of $(S_x-x\log x)/x$ tends to $e^{B(t)}$ as
$x\to\infty$, for every real $t$. The limit is continuous at $0$, so by
L\'evy's continuity theorem it is the characteristic function of a
probability law. It is integrable, so this law has the bounded continuous
density $f_L$ of Theorem~\ref{thm:landau}. For $0<c<1$ and real $t$ we have
$B(ct)=cB(t)-ict\log c$, and hence
\[
 e^{B(t)}=e^{B(ct)}\cdot e^{B((1-c)t)}\,e^{it[(1-c)\log(1-c)+c\log c]} .
\]
The last two factors form the characteristic function of a scaled and shifted
copy of the law. So the law is of class $L$ (self-decomposable), and $f_L$ is
unimodal by~\cite{yamazato1978}.

\emph{Proof of (a).} By Fourier inversion on $\Z$ and the substitution
$\theta=t/x$,
\[
 x\,\CP_N(m)=\frac1{2\pi}\int_{-\pi x}^{\pi x}e^{-i\lambda_mt}e^{A(t)}\,dt,\qquad
 A(t)=x\Psi_N(t/x)-it\log x .
\]
Let $\lvert t\rvert\le\pi x$. By \eqref{eq:bulk-AB} and Fact 1,
\[
 \lvert A(t)-B(t)\rvert\le\frac{t^2}x\bigl(0.631(\lvert\log\lvert t\rvert\rvert+\log x)+1.53\bigr)+\frac{2x}N,
\]
and by Facts 1 and 3,
$\max(\operatorname{Re}A(t),\operatorname{Re}B(t))\le-\frac{\lvert t\rvert}{2\pi}+\frac{2x}N$.
Now $\lvert e^A-e^B\rvert\le\lvert A-B\rvert\,e^{\max(\operatorname{Re}A,\operatorname{Re}B)}$,
and $\frac{2x}N\le\frac12$. The integrals of $t^2e^{-\lvert t\rvert/2\pi}$,
of $t^2\lvert\log\lvert t\rvert\rvert e^{-\lvert t\rvert/2\pi}$ and of
$e^{-\lvert t\rvert/2\pi}$ over the real line are finite. So
\[
 \frac1{2\pi}\int_{-\pi x}^{\pi x}\lvert e^{A(t)}-e^{B(t)}\rvert\,dt
 \le C\Bigl(\frac{1+\log x}x+\frac xN\Bigr).
\]
Beyond $\pi x$ we use $\lvert e^{B(t)}\rvert=e^{-\pi\lvert t\rvert/2}$, which
gives $\frac1{2\pi}\int_{\lvert t\rvert>\pi x}\lvert e^{B(t)}\rvert\,dt
=\frac2{\pi^2}e^{-\pi^2x/2}\le\frac Cx$. Since
$f_L(\lambda_m)=\frac1{2\pi}\int_{\mathbb R}e^{-i\lambda_mt}e^{B(t)}\,dt$,
this proves (a).

\emph{Proof of (b).} By Theorem~\ref{thm:cp},
$\lvert xf_N(m)-x\,\CP_N(m)\rvert\le x\delta_B$. Add (a).

\emph{Proof of (c).} Let $x\to\infty$ and $x^3\log^2N/N\to0$. Then
$x\delta_B\to0$ and $x/N\to0$, so the bound of (b) tends to $0$. Let $U$ be
uniform on $[0,1)$ and independent of $M_N$. Then $(M_N+U-x\log x)/x$ has
the density $g_N$ with $g_N(\lambda)=xf_N(m)$ for
$\lambda_m\le\lambda<\lambda_m+\frac1x$. On a compact set of $\lambda$ the
relevant $m$ are nonnegative once $x$ is large, and
$\lvert\lambda-\lambda_m\rvert<\frac1x$. As $f_L$ is uniformly continuous,
$g_N\to f_L$ uniformly on compact sets. By Scheff\'e's lemma $g_N\to f_L$ in
$L^1$, and since $U/x\to0$ this gives the convergence in law. At the
crossing $x_N\sim2\log N$ by Theorem~\ref{thm:erw-crossing}, so both
conditions hold. \qed

\subsection{Proof of Lemma~\ref{lem:landau-mode}}\label{app:bulk-mode}

\emph{Step 1: a real integral.} We claim that for every real $\lambda$ and
$k\ge0$,
\begin{equation}\label{eq:landau-real}
 f_L^{(k)}(\lambda)=\frac1\pi\int_0^\infty(-s)^k\,e^{-s\log s-\lambda s}\sin(\pi s)\,ds .
\end{equation}
Since $B(-t)=\overline{B(t)}$, the definition of $f_L$ gives
$f_L(\lambda)=\frac1\pi\operatorname{Re}\int_0^\infty\Phi(-it)\,dt$ with
$\Phi(u)=\exp(u\log u+\lambda u)$. The function $\Phi$ is holomorphic off the
negative real axis. Put $u=\rho e^{i\varphi}$ with
$-\pi<\varphi\le-\frac\pi2$. Then
$\operatorname{Re}(u\log u)=\rho(\log\rho\cos\varphi-\varphi\sin\varphi)$,
where $\cos\varphi\le0$ and $\varphi\sin\varphi\ge0$. For
$\log\rho\ge\lvert\lambda\rvert+1$ this gives
$\operatorname{Re}(u\log u+\lambda u)\le-\rho(\lvert\cos\varphi\rvert+\varphi\sin\varphi)\le-\rho$,
because $\lvert\cos\varphi\rvert+\varphi\sin\varphi\ge1$ on
$[-\pi,-\frac\pi2]$ (for $\psi=\varphi+\pi\in[0,\frac\pi2]$ it equals
$\cos\psi+(\pi-\psi)\sin\psi\ge\cos\psi+\sin\psi\ge1$). Also $\Phi$ is bounded
near $0$. So by Cauchy's theorem we may turn the ray $u=-it$ to the lower
side of the negative real axis, where $u=-s$ and $\log u=\log s-i\pi$. There
$u\log u=-s\log s+i\pi s$, and $dt=i\,du=-i\,ds$. Hence
\[
 f_L(\lambda)=\frac1\pi\operatorname{Re}\Bigl[-i\int_0^\infty
 e^{-s\log s-\lambda s}\,e^{i\pi s}\,ds\Bigr]
 =\frac1\pi\int_0^\infty e^{-s\log s-\lambda s}\sin(\pi s)\,ds .
\]
The integrand decays faster than every exponential, so we may differentiate
under the integral sign. This proves \eqref{eq:landau-real}. Also $f_L''$
is bounded on the real line, because $t^2e^{B(t)}$ is integrable and we may
differentiate the Fourier formula for $f_L$ twice.

\emph{Step 2: certified enclosures.} Let $-0.223\le\lambda\le0$ and
$k\in\{1,2\}$, and split the integral in \eqref{eq:landau-real} at
$\delta=2^{-20}$ and at $T=30$. On $[0,\delta]$ we have $-s\log s\le1/e$,
$-\lambda s\le0.223\,s$ and $\lvert\sin\pi s\rvert\le\pi s$, so this part is
at most $\pi e^{1/e+0.223\delta}\delta^{k+2}/(k+2)$ in absolute value. On
$[T,\infty)$ we have $s\log s+\lambda s\ge cs$ with $c=\log30-0.223>3$, so
this part is at most
$\int_T^\infty s^2e^{-cs}\,ds\le T^2e^{-cT}$. Both are below $10^{-17}$. On
$[\delta,T]$ the integrand is holomorphic in a neighborhood of the path, and
we enclose the integral with the rigorous integration algorithm
of~\cite{johansson2018}, in ball arithmetic with 128 bits. The
parameter $\lambda$ enters as a ball, so one run encloses the integral for
every $\lambda$ in an interval. The script \texttt{certify\_elephant.py} does
this and reports
\[
 f_L'(-0.2229)\in[9.232\cdot10^{-6},9.234\cdot10^{-6}],\qquad
 f_L'(-0.2227)\in[-6.548\cdot10^{-6},-6.546\cdot10^{-6}],
\]
and $f_L''(\lambda)\in[-0.0792,-0.0786]$ for every
$\lambda\in[-0.2229,-0.2227]$. In the same way it finds that $f_L'$ is
positive at $-0.22278299$ and negative at $-0.22278297$.

\emph{Step 3: the mode.} Write $[a,b]=[-0.2229,-0.2227]$. By Step 2, $f_L'$
is strictly decreasing on $[a,b]$ with $f_L'(a)>0>f_L'(b)$. So $f_L'$ has
exactly one zero $\lambda_0$ in $[a,b]$, and $f_L$ is strictly increasing on
$[a,\lambda_0]$ and strictly decreasing on $[\lambda_0,b]$. Since $f_L$ is
unimodal, there is a point $\lambda_*$ such that $f_L$ is nondecreasing on
$(-\infty,\lambda_*]$ and nonincreasing on $[\lambda_*,\infty)$. If
$\lambda_*>\lambda_0$, then $f_L$ would be nondecreasing just to the right of
$\lambda_0$, which is false. In the same way $\lambda_*<\lambda_0$ is
impossible. So $\lambda_*=\lambda_0$. By Step 2 again,
$\lambda_0\in(-0.22278299,-0.22278297)$. This proves (a) and (b).

\emph{Step 4: the constant $\kappa_L$.} Let
$\delta_0=\min(\lambda_0-a,b-\lambda_0)>0$. For
$\lvert\lambda-\lambda_0\rvert\le\delta_0$, Taylor's formula with
$f_L'(\lambda_0)=0$ and $f_L''\le-0.078$ on $[a,b]$ gives
$f_L(\lambda_0)-f_L(\lambda)\ge0.039(\lambda-\lambda_0)^2$. For
$\lvert\lambda-\lambda_0\rvert>\delta_0$ the monotonicity of Step 3 gives
$f_L(\lambda)\le f_L(\lambda_0\pm\delta_0)\le f_L(\lambda_0)-0.039\,\delta_0^2$.
Since $\delta_0<1$, (c) holds with $\kappa_L=0.039\,\delta_0^2$. \qed

\subsection{Proof of Corollary~\ref{cor:mode-location}}\label{app:bulk-cor}

The method is to compare $f_N$ at a mode with $f_N$ at the integer $m'$
nearest to the predicted place $x\log x+\lambda_0x$. By
Theorem~\ref{thm:landau}(b) both values are close to values of $f_L$, and
$f_L$ falls off quadratically away from $\lambda_0$. So a mode cannot be far
from $m'$.

We first prove a statement at finite $N$. Assume the standing assumptions of
Section~\ref{sec:shape}, let $\kappa_L$ be as in
Lemma~\ref{lem:landau-mode}(c), and let $E_1$ be the bound of
Theorem~\ref{thm:landau}(b). Put
\[
 E'=E_1+\frac{16x^2H_N}{N^2}+\frac{\lVert f_L''\rVert_\infty}{8x^2}.
\]
We claim that if $E'<\kappa_L/2$ and $m_h\le N/4$, then every mode $m$ of
$f_N$ satisfies
$\lvert m-x\log x-\lambda_0x\rvert\le x\sqrt{2E'/\kappa_L}$, and so does
$m^*$ for every mode $a^*\ge N/2$ of $P_N$.

Let $m'$ be the integer nearest to $x\log x+\lambda_0x$. This number
increases in $x\ge1$ and is at least $\lambda_0>-\frac12$, so $m'\ge0$. Also
$\lvert\lambda_{m'}-\lambda_0\rvert\le\frac1{2x}$, and since
$f_L'(\lambda_0)=0$,
$f_L(\lambda_{m'})\ge f_L(\lambda_0)-\lVert f_L''\rVert_\infty/(8x^2)$. Let
$\tilde m$ be a mode of $f_N$ or $\tilde m=m^*$. We claim that
$f_N(\tilde m)\ge f_N(m')-16xH_N/N^2$. For a mode of $f_N$ this is clear.
For $m^*$ we have
$f_N(m^*)+f_N(N-m^*)=2P_N(N-m^*)\ge2P_N(N-m')\ge f_N(m')$. Here
$N-m^*\ge\lceil N/2\rceil\ge\lceil N/4\rceil\ge m_h$. So
$f_N$ is nonincreasing on $\{\lceil N/4\rceil,\dots,N\}$ by
Lemma~\ref{lem:unimodal}. There are at least $N/4$ integers from
$\lceil N/4\rceil$ to $\lceil N/2\rceil$. So Lemma~\ref{lem:B-sandwich}
(with $a=\lceil N/4\rceil\ge m_h$ and $b=\lceil N/2\rceil$) gives the second
inequality below, and Lemma~\ref{lem:B-tail} gives the third,
\[
 f_N(N-m^*)\le f_N(\lceil N/2\rceil)\le\frac{\Prob(M_N\ge\lceil N/4\rceil)}{N/4}
 \le\frac{16xH_N}{N^2}.
\]
Now Theorem~\ref{thm:landau}(b), used at $\tilde m$ and at $m'$, gives
\[
 f_L(\lambda_{\tilde m})\ge xf_N(\tilde m)-E_1\ge xf_N(m')-\frac{16x^2H_N}{N^2}-E_1
 \ge f_L(\lambda_{m'})-2E_1-\frac{16x^2H_N}{N^2}\ge f_L(\lambda_0)-2E'.
\]
Since $2E'<\kappa_L$, Lemma~\ref{lem:landau-mode}(c) gives
$(\lambda_{\tilde m}-\lambda_0)^2\le2E'/\kappa_L$, which is the claim.

Now take $x=x_N$ at the crossing. By Theorem~\ref{thm:erw-crossing},
$x\sim2\log N$. So $x\delta_B$, $x/N$, $16x^2H_N/N^2$ and
$\lVert f_L''\rVert_\infty/(8x^2)$ are all $O(\log x/x)$, and
$E'=O(\log x/x)$. For large $N$ the standing assumptions hold, $E'<\kappa_L/2$,
and Lemma~\ref{lem:B-mode} gives $m_h\le A_0x\log^2(x+1)\le N/4$. Hence
$\lvert m^*-x\log x-\lambda_0x\rvert\le x\sqrt{2E'/\kappa_L}=O(\sqrt{x\log x})$,
and the same holds for $m_h$. \qed
